\ifluatex\else
% textcomp defines 00B1 wrong, so fix it here
\DeclareUnicodeCharacter{2192}{\ensuremath{\rightarrow}}
\DeclareUnicodeCharacter{2194}{\ensuremath{\leftrightarrow}}
\DeclareUnicodeCharacter{2212}{\ensuremath{-}}
\DeclareUnicodeCharacter{22EE}{\ensuremath{\vdots}}
\DeclareUnicodeCharacter{22EF}{\ensuremath{\cdots}}
\DeclareUnicodeCharacter{00B1}{\ensuremath{\pm}}
\DeclareUnicodeCharacter{00D7}{\ensuremath{\times}}
\DeclareUnicodeCharacter{221E}{\ensuremath{\infty}}
\DeclareUnicodeCharacter{2248}{\ensuremath{\approx}}
\DeclareUnicodeCharacter{2264}{\ensuremath{\leq}}
\DeclareUnicodeCharacter{2265}{\ensuremath{\geq}}
\DeclareUnicodeCharacter{03BC}{\micro}
\DeclareUnicodeCharacter{2206}{\ensuremath{\Delta}}
\fi

% KU Leuven colour definitions
\makeatletter
\@ifpackageloaded{xcolor}{%
    \definecolor{kuldefault}{HTML}{00407a}%
    \definecolor{kulbright}{HTML}{52bdec}%
    \definecolor{kulleft}{HTML}{1d8db0}%
    \definecolor{kulright}{HTML}{116e8a}%
    \definecolor{kuldark}{HTML}{123C75}% heading background
    \definecolor{kullight}{HTML}{F0F7FA}% block background
    \definecolor{kulbg}{RGB}{29,141,176}% background fill
    \definecolor{kulyellow}{HTML}{BC8F00}%
    \definecolor{kulorange}{HTML}{BC6E00}%
    \definecolor{kulgreen}{HTML}{007F4F}%
    \definecolor{kulred}{HTML}{FF4422}%
}{}
\makeatother

% Some handy abbreviations
\newcommand{\eg}{e.\,g.\xspace}
\newcommand{\ie}{i.\,e.\xspace}

% SI units using siunitx
\DeclareSIUnit{\octave}{octave}
\DeclareSIUnit{\pbdrunit}{\%BDR}
\DeclareSIUnit{\plbdrunit}{\%LBDR}
\DeclareSIUnit{\currentunit}{cu}
\DeclareSIUnit{\pulsespersecondunit}{pps}

% Convenience commands (assuming you want to keep the same syntax)
\newcommand{\pps}[1]{\SI{#1}{\pulsespersecondunit}}
\newcommand{\hz}[1]{\SI{#1}{\hertz}}
\newcommand{\khz}[1]{\SI{#1}{\kilo\hertz}}
\newcommand{\mhz}[1]{\SI{#1}{\mega\hertz}}
\newcommand{\db}[1]{\SI{#1}{\deci\bel}}
\newcommand{\us}[1]{\SI{#1}{\micro\second}}
\newcommand{\ms}[1]{\SI{#1}{\milli\second}}
\newcommand{\s}[1]{\SI{#1}{\second}}
\newcommand{\pc}[1]{\SI{#1}{\percent}}
\newcommand{\pmil}[1]{\SI{#1}{\per\thousand}}
\newcommand{\pcdb}[1]{\SI{#1}{\percent\per\deci\bel}}
\newcommand{\de}[1]{\SI{#1}{\degree}}
\newcommand{\uv}[1]{\SI{#1}{\micro\volt}}
\newcommand{\nv}[1]{\SI{#1}{\nano\volt}}
\newcommand{\bdr}[1]{\SI{#1}{\pbdrunit}}
\newcommand{\lbdr}[1]{\SI{#1}{\plbdrunit}}
\newcommand{\cu}[1]{\SI{#1}{\currentunit}}
\newcommand{\nhl}[1]{\SI{#1}{\deci\bel\textnormal{nHL}}}
\newcommand{\hl}[1]{\SI{#1}{\deci\bel\textnormal{HL}}}
\newcommand{\spl}[1]{\SI{#1}{\deci\bel\textnormal{SPL}}}
\newcommand{\spla}[1]{\SI{#1}{\deci\bel\textnormal{SPL(A)}}}
\newcommand{\dbfs}[1]{\SI{#1}{\deci\bel\textnormal{FS}}}
\newcommand{\pespl}[1]{\SI{#1}{\deci\bel\textnormal{peSPL}}}
\newcommand{\pespla}[1]{\SI{#1}{\deci\bel\textnormal{peSPL(A)}}}
\newcommand{\sd}[1]{\ensuremath{\textnormal{SD}=#1}}
\newcommand{\se}[1]{\ensuremath{\textnormal{SE}=#1}}
\newcommand{\ir}[1]{\ensuremath{\textnormal{IR}=#1}}
\newcommand{\el}[2]{\ensuremath{\textnormal{#1}_{\textnormal{#2}}}}

% pandoc compatibility
\providecommand{\tightlist}{\setlength{\itemsep}{0pt}\setlength{\parskip}{0pt}}

\newcommand{\deflen}[2]{%
    \expandafter\newlength\csname #1\endcsname
    \expandafter\setlength\csname #1\endcsname{#2}%
}
